%% 通用示例模板：以下评论均为虚构示例，使用时替换为用户的真实审稿原文。
%% 保留原评论编号。下方表格表示拟修改内容，不表示工作已经完成。

%% 尊敬的编辑和审稿人：
Dear Editor and Reviewers:

%% [开头说明占位符：在此填写论文题目和稿件编号。]
[Opening statement, manuscript title, and manuscript ID to be completed.]

%% 此致，[作者署名待填写。]
Yours sincerely,\par
[Author names]

%% 表格按拟修改章节关联原评论；同一行可对应多条评论。
\begingroup
\footnotesize
\renewcommand{\arraystretch}{1.15}
\setlength{\tabcolsep}{3.5pt}
\setlength{\LTpre}{5pt}
\setlength{\LTpost}{6pt}
\begin{longtable}{
>{\raggedright\arraybackslash}p{2.6cm}
>{\raggedright\arraybackslash}p{3.3cm}
>{\raggedright\arraybackslash}p{9.0cm}}
%% 拟修改内容汇总（占位符）。
\caption{Summary of Planned Revisions (Placeholders)}
\label{tab:revision_summary}\\
\toprule
%% 论文章节；关联评论；拟修改内容。
\textbf{Manuscript Section} & \textbf{Related Comments} & \textbf{Planned Revision} \\
\midrule
\endfirsthead
\multicolumn{3}{c}{\textit{Table~\thetable\ continued from the previous page}}\\
\toprule
\textbf{Manuscript Section} & \textbf{Related Comments} & \textbf{Planned Revision} \\
\midrule
\endhead
\midrule
\multicolumn{3}{r}{\textit{Continued on the next page}}\\
\endfoot
\bottomrule
\endlastfoot
{[Section to be verified]} & R1-C1 & [Proposed clarification; not yet completed.]\\
\midrule
{[Section to be verified]} & R1-C2; R2-C1 & [Shared analysis to be planned; results pending.]\\
\end{longtable}
\endgroup

\clearpage
%% 对审稿人 1 的回复。
\section*{Response to Reviewer 1}

%% 评论 1。
\subsection*{Comment 1}
\begin{cmt}{}{}
%% 中文忠实翻译：
%% 请定义目标总体，并说明为什么所选评估指标适用于该总体。
Please define the target population and explain why the chosen evaluation metric is appropriate for that population.
\end{cmt}

%% 建议回复逻辑（供作者参考）：
%% 核对总体与纳入标准 → 将评估指标对应到研究目标 → 据此确定需要补充的定义与解释。
\responseheading{1}

%% 原子问题 R1-C1-a：目标总体如何定义？
%% 小标题中文：目标总体。
\textbf{(a) Target Population}\par
\responseplaceholder

%% 原子问题 R1-C1-b：为什么所选评估指标适用于该总体？
%% 小标题中文：评估指标的适用性。
\textbf{(b) Suitability of the Evaluation Metric}\par
\responseplaceholder

%% 评论 2。
\subsection*{Comment 2}
\begin{cmt}{}{}
%% 中文忠实翻译：
%% 仅凭总体结果，无法判断该方法对历史记录有限的用户是否也表现良好。
%% 请提供按历史记录长度划分的评估。
The aggregate results do not establish whether the method also performs well for users with limited history. Please provide an evaluation by history length.
\end{cmt}

%% 建议回复逻辑（供作者参考）：
%% 核对是否已有历史长度分组结果 → 如有缺口则规划分组评估 → 待结果核实后界定适用范围；关联 R2-C1。
\responseheading{2}

%% 原子问题 R1-C2-a：不同历史记录长度下的方法表现如何？共享组 G1，关联 R2-C1-a。
%% 小标题中文：按历史记录长度划分的表现。
\textbf{(a) Performance by History Length}\par
\responseplaceholder

\clearpage
%% 对审稿人 2 的回复。
\section*{Response to Reviewer 2}

%% 评论 1。
\subsection*{Comment 1}
\begin{cmt}{}{}
%% 中文忠实翻译：
%% 该方法对历史记录较短的用户可能效果较差。你们能否按历史记录长度报告结果？
The method may be less effective for users with short histories. Could you report results by history length?
\end{cmt}

%% 建议回复逻辑（供作者参考）：
%% 复用 R1-C2 的历史长度分组评估 → 根据核实结果回应潜在表现差异 → 统一相关适用范围表述。
\responseheading{1}

%% 原子问题 R2-C1-a：能否用分组结果评估短历史用户的表现？共享组 G1，关联 R1-C2-a。
%% 小标题中文：短历史用户的表现。
\textbf{(a) Performance for Users with Short Histories}\par
\responseplaceholder
